\documentclass{article}

\usepackage{amsmath,amssymb}
\usepackage{xurl}
\usepackage[hidelinks]{hyperref}
\setlength{\emergencystretch}{2em}

% Shared commands for notes in docs/paper-gaps/.

\DeclareUnicodeCharacter{2080}{\ifmmode{}_{0}\else\textsubscript{0}\fi}
\DeclareUnicodeCharacter{2081}{\ifmmode{}_{1}\else\textsubscript{1}\fi}
\DeclareUnicodeCharacter{2082}{\ifmmode{}_{2}\else\textsubscript{2}\fi}
\DeclareUnicodeCharacter{2099}{\ifmmode{}_{n}\else\textsubscript{n}\fi}

\newcommand{\Lean}{\textsc{Lean}}
\ExplSyntaxOn
\NewDocumentCommand{\leanid}{m}{
  \ifmmode
    \textsf{\detokenize{#1}}
  \else
    \group_begin:
    \urlstyle{sf}
    \tl_set:Nn \l_tmpa_tl {#1}
    \tl_replace_all:Nnn \l_tmpa_tl {\_} {_}
    \exp_args:NV \path \l_tmpa_tl
    \group_end:
  \fi
}
\ExplSyntaxOff
\providecommand{\tr}{\operatorname{tr}}
\newcommand{\tnleanIssueBase}{https://github.com/LionSR/TNLean/issues/}
\newcommand{\tnleanPrBase}{https://github.com/LionSR/TNLean/pull/}
\newcommand{\tnleanDocBase}{https://sirui-lu.com/TNLean/docs/}
\newcommand{\ghissue}[1]{\href{\tnleanIssueBase#1}{GitHub issue \##1}}
\newcommand{\ghpr}[1]{\href{\tnleanPrBase#1}{PR \##1}}
\newcommand{\doclink}[1]{\href{\tnleanDocBase#1.html}{\path{#1}}}

% Dirac notation and the identity operator, as in blueprint/src/macros/common.tex.
% \providecommand keeps note-local definitions authoritative.
\usepackage{dsfont}
\providecommand{\ket}[1]{|#1\rangle}
\providecommand{\bra}[1]{\langle #1|}
\providecommand{\braket}[2]{\langle #1 | #2 \rangle}
\providecommand{\ketbra}[2]{|#1\rangle\!\langle #2|}
\providecommand{\Id}{\mathds{1}}
\providecommand{\one}{\mathds{1}}
\providecommand{\C}{\mathbb{C}}
\providecommand{\MN}[1]{M_{#1}(\C)}
\providecommand{\MD}{M_D(\C)}

% External referencing for paper sources.  The build helper writes sanitized
% auxiliary files under build/paper-aux/ and excludes duplicated source labels.
\usepackage{xr}

% Import a generated paper auxiliary file with a caller-supplied label prefix.
% The candidate paths support builds from docs/paper-gaps/, the repository root,
% and build/paper-aux/ itself.
\newcommand{\importpaperaux}[2]{%
  \IfFileExists{../../build/paper-aux/#2.aux}{%
    \externaldocument[#1]{../../build/paper-aux/#2}%
  }{%
    \IfFileExists{build/paper-aux/#2.aux}{%
      \externaldocument[#1]{build/paper-aux/#2}%
    }{%
      \IfFileExists{#2.aux}{%
        \externaldocument[#1]{#2}%
      }{}%
    }%
  }%
}

% General external reference macros
\newcommand{\paperref}[2]{\ref{#1-#2}}
\newcommand{\papereqref}[2]{(\ref{#1-#2})}

% CPSV16 source (1606.00608 / MPDO-22-12-17-2.tex) external document and shortcuts
\importpaperaux{cpsv16-}{1606.00608/MPDO-22-12-17-2-tnlean}
\newcommand{\cpsvref}[1]{\paperref{cpsv16}{#1}}
\newcommand{\cpsveqref}[1]{\papereqref{cpsv16}{#1}}

% Verdict marker: \gapnote{<kind>}{<status>}. Kind is the policy
% classification (clarification, local-correction, scope-restriction,
% unfaithful, false-source, open-gap); status is open, wip (elimination
% underway), resolved, or historical. Machine-read by the site index;
% prints a verdict line.
\newcommand{\gapnote}[2]{\par\noindent\textbf{Verdict:} #1 (#2).\par}


\title{The Trace Sign in Wolf's Lorentz-Cone Converse}
\author{}
\date{2026-06-23}

\begin{document}
\maketitle

\section{The assertion}

Wolf's Proposition~3.9 states two elementary trace inequalities for Hermitian
matrices.\footnote{\cite[Proposition~3.9]{Wolf2012Quantum};
local source \path{Notes/WolfNoteTexSource/ch03_positive_not_completely.tex},
lines 693--722.}  The forward implication says that a positive semidefinite
matrix satisfies
\begin{align}
  \tr(A^2) \le \tr(A)^2.
\end{align}
The converse is printed as
\begin{align}
  \tr(A)^2 \ge (d-1)\tr(A^2) \Longrightarrow A \ge 0.
  \tag{\path{Wolf Prop. 3.9 converse}}
\end{align}

\section{The point at issue}

The printed converse cannot be true without a sign condition on the trace.  For
example, if \(A=-I_d\) and \(d\geq 1\), then \(A\) is Hermitian and not positive
semidefinite, while
\begin{align}
  \tr(A)^2=d^2 \ge d(d-1)=(d-1)\tr(A^2).
\end{align}
Thus the squared inequality only controls the Lorentz cone up to time
orientation; it does not distinguish the positive and negative trace sheets.

This is also visible in the proof printed after the proposition.  The argument
decomposes \(A=A_+-A_-\) and uses the estimate
\begin{align}
  \tr(A) \le \tr(A_+).
\end{align}
To contradict the squared inequality one needs the left hand side to have the
same orientation as the nonnegative quantity being bounded.  A nonnegative
trace hypothesis supplies exactly this orientation.

\section{Corrected statement}

The formal theorem proves the trace-nonnegative version:
\begin{align}
  A=A^\dagger,\qquad
  \tr(A)\ge 0,\qquad
  (d-1)\tr(A^2)\le \tr(A)^2
  \Longrightarrow A\ge 0.
\end{align}
In the complex matrix formalization the traces are written with real parts,
which agrees with the displayed real quantities for Hermitian \(A\).\footnote{
Formal declaration
\leanid{Matrix.IsHermitian.posSemidef_of_trace_re_nonneg_of_card_sub_one_mul_trace_sq_re_le}
in \doclink{TNLean/Analysis/MatrixTraceInequalities}; tracked by
\ghissue{3402}.}

The proof is a finite eigenvalue argument.  If some eigenvalue \(\lambda_0\) is
negative, write
\begin{align}
  T=\sum_i\lambda_i,\qquad
  R=\sum_{i\ne 0}\lambda_i .
\end{align}
Since \(T\ge0\) and \(T=\lambda_0+R\), one has \(T<R\), and hence
\(T^2<R^2\).  Cauchy's inequality for the remaining \(d-1\) eigenvalues gives
\begin{align}
  R^2\le (d-1)\sum_{i\ne 0}\lambda_i^2
       \le (d-1)\sum_i\lambda_i^2,
\end{align}
contradicting the assumed Lorentz inequality.  Therefore all eigenvalues are
nonnegative.

\section{Conclusion}

The unrestricted printed converse should not be marked as formalized.  The
formal development records the future-cone, trace-nonnegative converse as the
mathematically valid statement used in the positive-map discussion.

\bibliographystyle{alpha}
\bibliography{references}

\end{document}
